\documentclass[10pt]{article}

\usepackage[utf8]{inputenc}

\usepackage[T1]{fontenc}

\usepackage{amsmath}

\usepackage{amsfonts}

\usepackage{amssymb}

\usepackage[version=4]{mhchem}

\usepackage{stmaryrd}



\begin{document}

где



\[

\hat{h}_{\mu}^{j}=h_{\mu}^{j}-\min _{i=1, \ldots, h} h_{\mu}^{i},

\]



переводит пространство параметров на всю границу $n$-мерного положительного октанта. Это условие выполняется для семейства потенциалов «общего положения». Т е о р е ма П и рогова - Си ная (см. [1]-[3]) вместе с ее уточнением (см. [4]) утверждает, что существуют константы $\beta_{0}$, $\varepsilon_{0}>0$ такие, что при всех $\beta \geqslant \beta_{0}$ существует гомеоморфизм $\tau_{\beta}$ окрестности нулевой точки границы $n$-мерного положительного октанта $o^{n}$ в пространстве параметров $\mu$ такой, что если точка этой границы $x \in \partial o^{n}$ имеет $q$ нулевых координат, то для потенциала $U_{\mu} \mathrm{c} \mu=\Re_{\beta}(x)$ существует ровно $q$ различных регулярных периодич. состояний с обратной температурой $\beta$.



В случае, когда потенциал $U$ обладает нек-рыми симметриями, то аналогичной симметрией обладают и гиббсовские состояния с потенциалами $U_{\mu}$. Для такой ситуации результат П.- С.т. был известен и ранее (см. [5], [6]). Известны также обобщения П.- С.т. на более сложные ситуации (см.[7]-[9]).



Лит.: [1] П и р о го в С. А., С и н а й Я. Г., «Функциональный анализ и его прилож.», 1974, т. 8, № 1, с. 25-31; [2] и х же, «Теоретич. и математич. физика", 1975, т. 25, № 3, с. 358-69; [3] С и на й Я. Г., Теория фазовых переходов, М., 1980; [4] Zahradnik M., «Сomm. Math. Phys.», 1984, v. 93, p. 559-81; [5] Герцик В. М., Добруш и н Р. Л., «Функциональный анализ и его прилож.», 1974, т. 8, в. 3, с. 12-25; [6] Ге рц и к В. М., «Известия АН СССР. Сер. математическая», 1976, т. 40, с. 448-62; [7] Д и на б у рг Е. И., С инай Я. Г., «Труды конференции по математич. физике», Дубна, 1984; [8] I mbrie J. Z., "Comm. Math. Phys.», 1981, v. 82, p. 261-304; [9] Dobrushin R. L., Zahradnik M., B KH.: Mathematical problems of statistical mechanics and dynamics, Dordrecht - Boston, 1986, p. 1-123.



ПИРСИ ИНТЕГРАЛ - интеграл вида



\[

\operatorname{Pe}\left(\lambda_{1}, \lambda_{2}\right)=\int_{-\infty}^{\infty} e^{i\left(x^{4} / 8-\lambda_{1} x^{2} / 2+\lambda_{2} x\right)} d x

\]



Фаза П.и.- версальная деформация критич. точки $x^{4} / 8$ (типа $\left.A_{3}\right)$. Функция Pe $\left(\lambda_{1}, \lambda_{2}\right)$ называется функци е й П и рс и и возникает при исследовании волнового поля вблизи точки возврата каустики. С. В. Чмутов. ПИРСОНА УРАВНЕНИЕ - см. Ортогональные многочлены. ПЛАВЛЕНИЯ КРИВАЯ - см. Фазовая диаграмма.



ПЛАНКА ЗАКОН ИЗЛУЧЕНИЯ, Ф о р м У ла Пл а н Ка,закон распределения энергии в спектре равновесного излучения при определенной температуре $T$. Был впервые выведен М. Планком (M. Planck, 1900) на основе гипотезы квантования энергии вещества. М. Планк моделировал вещество совокупностями гармонич. осцилляторов различной частоты $v$ - резонаторов, испускающих и поглощающих излучение соответствующей частоты. Он предположил, что энергия вещества распределяется по резонаторам каждой частоты $v$ в виде дискретных порций $h v-$ квантов энергии ( $h$ - постоянная Планка). А. Эйнштейн (А. Einstein, 1916) вывел П. з. и. путем рассмотрения квантовых переходов для атомов, находящихся в равновесии с излучением. П. з. и. является частным случаем Бозе - Эйнштейна распределения.



П. з. и. дает спектральную зависимость (зависимость от частоты $v$ или длины волны $\lambda=c / v$ ) объемной плотности излучения (энергии излучения в единице объема) и пропорциональной ей испускательной способности абсолютно черного тела $\varepsilon=c u / 4$ (энергии излучения, испускаемой единицей его поверхности за единицу времени). Функции $u_{v}, T$ и $\varepsilon_{v, T}$ (или $u_{\lambda, T}$ и $\varepsilon_{\lambda, T}$ ), отнесенные к единице интервала частот (или длин волн), являются универсальными функциями от $v$ (или $\lambda$ ) и $T$, не зависящими от природы вещества, с к-рым излучение находится в равновесии.



П. з. и. выражается формулой



\[

u_{\mathrm{v}, T}=\frac{4}{c} \varepsilon_{\mathrm{v}, T}=\frac{8 \pi h v^{3}}{c^{3}} \frac{1}{\exp (h v / k T)-1}

\]



или



\[

u_{\lambda, T}=\frac{4}{c} \varepsilon_{\lambda, T}=\frac{8 \pi h c}{\lambda^{5}} \frac{1}{\exp (h c / \lambda k T)-1} .

\]



Максимум функции $\left({ }^{*}\right)$ с ростом $T$ смещается в сторону малых $\lambda$.



Из П. з. и. вытекают другие законы равновесного излучения. Интегрирование по $v$ (или $\lambda$ ) от 0 до $\infty$ дает значения полной объемной плотности излучения всех частот - за кон излучения Стефана - Больцмана:



\[

u_{T}=\int_{0}^{\infty} u_{v, T} d v=a T^{4}, \text { где } a=\frac{8 \pi^{5} h^{4}}{15 c^{3} h^{3}},

\]



и полной испускательной способности черного тела



\[

\varepsilon_{T}=\int_{0}^{\infty} \varepsilon_{v, T} d v=\sigma T^{4}, \text { где } \sigma=\frac{2 \pi^{5} h^{4}}{15 c^{2} h^{3}} .

\]



В области больших частот, когда энергия фотона много больше тепловой энергии $(h v \gg k T)$, П. з. и. переходит в закон излучения Вина:



\[

u_{v}, T^{=}\left(8 \pi h v^{3} / c^{3}\right) \exp (-h v / k T),

\]



в области малых частот $(h v \ll k T)-$ в закон излуче ния Рэлея - Джинса:



\[

u_{v, T}=\left(8 \pi v^{2} / c^{3}\right) k T \text {. }

\]



Таким образом, эти законы представляют собой предельные случаи П. з. и.



П. з. и. находится в согласии с экспериментальными данными, применяя его можно по этим данным вычислить значения $h$ и $k$. С помощью методов оптич. пирометрии можно на основе П. з. и. определять температуру нагретых тел.



Лит.: [1] Б ор н М., Атомная физика, пер. с англ., 3 изд., М., 1970; [2] С и в у и н Д. В., Общий курс физики, 2 изд., [т. 4] Oптика, M., 1985; [3] K u h n T. S., Klack-body theory and the quantum discontinuity, Oxf., 1978.



М. А. Ельяшевич.



ПЛАНКА МАССА - масса, определяемая соотношением



\[

M_{P l}=\sqrt{\frac{\hbar c}{G}}=1,22 \cdot 10^{19} \mathrm{Gev} / c^{2}=2,17 \cdot 10^{17}{ }_{\Gamma} \text {, }

\]



где $G$ - гравитационная постоянная, $\hbar$ - постоянная Планка, $c$ - скорость света. П. м.- одна из фундаментальных физич. величин. Ее значение в физике определяется тем, что это единственная величина размерности массы, к-рую можно составить из фундаментальных констант, описывающих три главных области физики: теорию тяготения $(G)$, теорию относительности $(c)$ и квантовую механику $(\hbar)$. Поэтому следует ожидать, что, во-первых, при описании взаимодействия элементарных частиц с энергиями, превышающими $M_{P l} c^{2}$, нельзя пренебрегать гравитационными эффектами и, во-вторых, что при таких энергиях гравитационные взаимодействия нельзя описывать классич. теорией, а нужно использовать квантовую теорию гравитации. Последовательное построение последней - одна из важнейших задач современной теоретич. физики. Одним из шагов в этом направлении является теория суперструн.



Jum.: [1] Proceedings of the fourth seminar on quantum gravity, 1984, Мау 25-29, М., 1987. А. Ю. Иенатьев.



ПЛАНКА ФОРМУЛА - см. Планка закон излучения.



ПЛАНШЕРЕЛЯ ТЕОРЕМА - сМ. ФУрье преобразование.



ПЛАСТИЧНОСТИ МАТЕМАТИЧЕСКАЯ ТЕОРИЯ - ТеОрия деформируемого пластичного твердого тела, в которой исследуются задачи, состоящие в определении полей вектора перемещений $u(x, t)$ или вектора скоростей $v(x, t)$, тензора деформации $\varepsilon_{i j}(x, t)$ или скоростей деформации $v_{i j}(x, t)$ и тензора напряжений $\sigma_{i j}(x, t)$, которые возникают в теле, занимающем область $\Omega$ с границей $S$. под действием массовых $K(x, t)$ и поверхностных $F(x, t)$ сил при заданных начальных и граничных условиях, а такхе в определении тех нагрузок и процессов, при которых равновесие (движе-





\end{document}